% SOURCE: results/tables/heuristic_evaluation.md "Inherent limitations".
% Numbers: 451 / 35 conflicting vectors / 0.580 macro-recall, all from
% results/tables/p1_model_metrics.csv and the cached frame.
\section{Limitations and Future Work}
\label{sec:limitations}

\subsection{Evaluation without users}

No mother, health worker or clinician has seen any of these renderings. That
is the limitation the rest depend on, and it bounds what the paper is allowed
to conclude: we have a design and a formative critique of our own artifact,
not evidence that the design works.

Four claims sit downstream of it. Whether a health worker reads the card as
intended is assumed and not shown; the nearest studies find that health
workers can struggle with AI explanations and that the design of an
explanation changes how well they understand
it~\citep{okolo2024easy,solanokamaiko2024explorable}. Whether
withholding magnitude from the mother helps or harms her is an argument from
the model's uncertainty, and the icon-array paragraph of
Section~\ref{sec:tiers-render} makes it, but the
risk-communication literature it departs from is empirical and our departure
is not. The $0.15$ margin is a design constant with no experiment behind it,
as Section~\ref{sec:tiers} states. And role remains a proxy: we design
against what each reader lacks in modality, literacy and decision authority,
but recruit nobody to check the proxy holds, which is the substance of Suresh
et al.'s objection to role-based tiering~\citep{suresh2021beyond}.

All four resolve the same way, and Section~\ref{sec:future} says how.
Everything here was produced and inspected on a workstation: no field claim,
no claim about performance under real clinic conditions, and no claim that
these renderings have reached anyone by any delivery channel.

\subsection{What one public dataset can support}
\label{sec:lim-data}

The work uses the UCI Maternal Health Risk dataset and nothing else. It is
public, carries generic three-level risk labels, and names no clinical
condition, so neither do we. Nothing here is clinically validated and nothing
predicts an outcome for a real patient.

Two properties bound the vehicle. Macro-recall is $0.580$ after
deduplication, the leak-free figure (Section~\ref{sec:data-leakage}); and 35
feature vectors, covering 71 of the 451 records, carry more than one risk
label, an ambiguity no model can resolve from six features. That is a ceiling
on the data, not a shortfall in the classifier.

% Derivations, all from results/tables/:
%  35 = p1_cohort_summary.csv conflicting_label_feature_vectors; 451 = n_clean.
%  71 = rows of p1_oof_predictions.csv with conflict_vector == True.

One assumption inside the gate deserves naming. The margin rule thresholds a
difference between random-forest class probabilities, which are vote
fractions and are not calibrated. We treat $0.15$ as a design
constant on an ordering, not a statement about probability.

The gate also has only one side. It tests the margin at the high boundary and
not the low one (Section~\ref{sec:tiers}). Out of fold, 22 low-risk
predictions sit inside the margin and render an unqualified \textsc{green} to
the mother; 3 of them are true high-risk records (Table~\ref{tab:gate}).

The most serious error mode is a missed alarm. Out of fold, 14 true high-risk
records receive \textsc{green}: 3 inside the $0.15$ margin and 11 outside it,
with margins from $0.230$ to $0.880$. Neither the current gate nor its
symmetric extension would reach them.

% Derivations: 22 and 3 = p1_gate_summary.csv, section
%  predicted_low_margin_below_0.15, row ungated_green_near_ties (total, high
%  risk). 14 = p1_gate_summary.csv true_high_reaches GREEN. 3 / 11 split:
%  p1_oof_predictions.csv rows with true_label == high risk and band == GREEN
%  (14), of which margin < 0.15 gives 3 and margin >= 0.15 gives 11; the 11
%  have margin min 0.2299999 and max 0.8799999 (typeset 0.230, 0.880).

\subsection{Who the renderings still exclude}
\label{sec:excl}

Each tier assumes something about its reader.

The mother's visual carries no rendered text, so a deaf or hard-of-hearing
mother gets nothing from it; the spoken sentence is her only channel and she
cannot hear it. A rendered caption or an SMS companion would close this, but
it reopens the literacy assumption the tier exists to avoid, so it needs user
input rather than a unilateral design call. Colour is the lamp's primary cue.
Position, the lit-versus-unlit luminance step and the spoken line are real
redundancies, but a reader who is both colour-blind and deaf is left with
position alone, which assumes the traffic-light convention. The clinician view
needs chart literacy and English. The health worker's card is English-only;
every string is a fixed template constant, so localising it is substitution
deferred to deployment. All three tiers ship as rasterised images of text.

\subsection{The voice is a stand-in}
\label{sec:voice}

Audio is synthesised with a general-purpose engine reached over the network.
It is not tuned for Indian-language health speech, handles acronyms
unpredictably, and needs connectivity the setting this design describes may
not have. Deployment-grade Indian-language synthesis (AI4Bharat, Bhashini) is
the stated path. What we treat as the artifact is the Hindi source string, not
the audio file.

\subsection{Future work}
\label{sec:future}

Studies with mothers, community health workers and clinicians are the first
item and the one the others wait on; they are gated on ethics clearance, which
this preprint does not have and does not claim. Two changes to the rule
are specified enough to name now. The gate should extend to the low boundary,
derating an ungated near-tie \textsc{green} to \textsc{amber} as the
red-to-amber move already does. On these data it would move the 22 near-ties
of Section~\ref{sec:lim-data} to \textsc{amber}; we report that count, and the
extension is not implemented. The forest's vote fractions should also be
calibrated, after which the threshold is worth re-asking. A third direction is
an explorable version of the health worker's card, which the findings of
Solano-Kamaiko et al.~\citep{solanokamaiko2024explorable} suggest.
Sections~\ref{sec:excl} and~\ref{sec:voice} name their own paths.
